\documentclass[10pt]{article}
\usepackage{nwa-guide}
\begin{document}
\guidetitle{Scope \& Sequence Workshop: Teacher's Guide}{NWA Mathematics \quad 2026--27}

\textbf{Two sites.} The \textbf{public site}
(\href{https://ryankcampbell.github.io/nwa-math-proficiencies/}{ryankcampbell.github.io/nwa-math-proficiencies})
is what parents, deans and students see. The \textbf{Workshop} (link in Ryan's email) is the same site
with edit buttons. Everything you do there saves at once to the department draft, with your name on it.
Nothing goes public until Ryan publishes, so experiment freely.

\textbf{Signing in.} Click the Workshop link while signed in to Google with your
\textbf{@nwacademy.org} account. If Google asks which account, pick the school one. A message
saying the workshop is for the math department means you're not on the list yet: tell Ryan.
Personal Gmail getting in the way? Open the link in a private window and sign in with the school account only.

\begin{multicols}{2}
\section*{1. Your course (Round 1)}
\ui{Courses} $\to$ your course. Every proficiency has a permanent code (\texttt{IM3-14}). Under each one:
\begin{itemize}
\item \textbf{Thread dropdown}: the common idea it belongs to (e.g.\ ``Exponential \& logarithmic
  functions''). Wrong home? Pick another.
\item \ui{Edit wording} \quad \ui{Retire} (you no longer teach it; its code is never reused)
\item \ui{+ Add proficiency} at the top for anything missing (it gets the next code).
\end{itemize}
Switch the view to \ui{By strand} to see your course thread by thread.

\section*{2. The scope grid (Round 1)}
\ui{Scope \& Sequence} $\to$ a strand tab. Rows are threads, columns are courses.
\textbf{Click your course's cell} in each row and write one short line: what your students do with that
idea. Leave \emph{n/a} if you don't touch it; mark optional or extension work as such.
\textit{More than a bar napkin, less than a research paper.}
Thread name wrong, or a note needed? \ui{Rename} or \ui{Note} beside it. A whole idea missing?
\ui{+ Add a thread}.

\section*{3. Overlaps and questions (Round 2)}
\ui{Review} lists proficiencies in neighboring courses that look alike. With your partner, decide each:
\textbf{Spiral} (taught again, deeper: say how), \textbf{Redundant} (one course should change), or
\textbf{Not the same}. See a gap or a sequencing problem? \ui{+ Raise a question}; it goes on the
November agenda.

\columnbreak
\section*{4. Resources (any time)}
On any proficiency, \ui{Add resource}: a title, then \textbf{paste a link} or \textbf{upload a file}
(PDF, image, doc; under 20\,MB), and a kind (video, practice, notes, interactive).
Free, no student login, your own work or openly licensed. Never a current quiz or test: uploads become
public. To remove one, ask Ryan.

\section*{5. Summer and make-up practice}
\ui{Assign practice}: tick proficiencies, or choose \ui{Summer before \emph{(course)}} to load
everything the next course assumes. Add a title, note and due date. Then either
\begin{itemize}
\item \ui{Create link}: a link, a \textbf{QR code} and ready-made \textbf{email text}; or
\item \ui{Make a class code}: five characters (e.g.\ \texttt{K7QM4}), live within a minute.
\end{itemize}
Students open the public site and type the code on the home page (no login), or open the link. They
get a checklist with resources and \textbf{self-check problems} (new numbers every try; graphs, sliders,
place-a-point). The app remembers their codes. \emph{No student information is stored anywhere}, so
never put names in a title or note. Try the demo code \texttt{NWA26}.

\section*{6. Send Ryan your files}
Old \textbf{exams and quizzes with keys}, worksheets, review packets, and \textbf{outside resources you
like} (the link and the proficiency it fits). Exams are never posted: Claude uses them to see what we
assess and to write \emph{new} practice problems in that style, which you approve before they go public.

\section*{Good to know}
\begin{itemize}
\item See colleagues' latest edits: \ui{refresh} in the teal bar.
\item Two people editing the same thing: the later save wins, and every version stays in the change log.
\item Only Ryan publishes; courses, strands and the teacher list live in his Sheet.
\end{itemize}
\end{multicols}
\end{document}
